\documentclass[a4paper]{article}
\usepackage[margin=20mm]{geometry}
\usepackage{graphicx}
\usepackage{float}
\renewcommand{\figurename}{Supplementary Fig.}
\renewcommand{\thefigure}{S\arabic{figure}}
\renewcommand{\tablename}{Supplementary Table}
\renewcommand{\thetable}{S\arabic{table}}

\begin{document}
\begin{center}
{\Large\textbf{Online Resource 1: Supplementary analyses for RQ1}\par}
\vspace{4mm}
{\small\textit{Cache Configuration Smells in GitHub Actions: A Large-Scale Empirical Study}\par}
{\small Empirical Software Engineering\par}
{\small Siyu Wang, Yang Zhang, Changrong Xie, Jiacheng Jie, Youren Chen, Zezhou Tang, Tanghaoran Zhang, Xinjun Mao\par}
{\small College of Computer Science and Technology, National University of Defense and Technology, Changsha, China\par}
{\small Corresponding author: Xinjun Mao (\texttt{xjmao@nudt.edu.cn})\par}
\end{center}

\section*{Workflow-level CCS co-occurrence}

Supplementary Table~S1 lists the eight most frequent unique pairs among the
11 operationalized CCS rules. Pair counts are workflows triggering both
rules; corpus share divides that count by all 118,402 workflows with usable detector results.
Jaccard divides the pair count by the number of workflows triggering at least
one rule in the pair. These summaries are descriptive; they do not establish
independent associations or causal relations. The complete pairwise matrix is
provided in \texttt{rq1\_cooccurrence.csv}.

\begin{table}[H]
\centering
\small
\caption{Selected workflow-level co-occurrences among the 11 CCS rules.}
\label{tab:supp-smell-cooccurrence}
\begin{tabular}{p{0.48\textwidth}rrr}
\hline
CCS pair & Workflows & Corpus share & Jaccard\\
\hline
No fallback cache + static cache key & 1,596 & 1.3480\% & 0.271382\\
No fallback cache + redundant cache update & 1,050 & 0.8868\% & 0.157256\\
No dependency cache + no fallback cache & 858 & 0.7246\% & 0.040785\\
No dependency cache + no docker layer cache & 362 & 0.3057\% & 0.018443\\
Redundant cache update + static cache key & 351 & 0.2964\% & 0.105121\\
No fallback cache + overlapping cache mechanisms & 314 & 0.2652\% & 0.051450\\
No dependency cache + redundant cache update & 313 & 0.2644\% & 0.017589\\
No fallback cache + overly specific cache key & 311 & 0.2627\% & 0.050925\\
\hline
\end{tabular}
\end{table}

\section*{NDC within supported dependency-installation opportunities}

Supplementary Table~S2 separates the 28,450 workflows with supported
dependency-installation metadata into mutually exclusive groups. The npm-compatible group covers npm, Yarn, and pnpm; the
Python-compatible group covers pip, Pipenv, Poetry, and uv. A workflow using
both groups appears only in the mixed row. These groups describe operations
observed in workflow jobs rather than repository primary language. Unsupported
dependency ecosystems and opaque dependency acquisition are outside the NDC
opportunity denominator. This metadata-based denominator includes operations excluded by the detector as auxiliary or cache warming. The rates are descriptive and do not establish an
ecosystem effect.

\begin{table}[H]
\centering
\small
\caption{NDC detection among workflows with supported dependency-installation opportunities.}
\label{tab:supp-ndc-opportunity}
\begin{tabular}{lrrr}
\hline
Dependency operations & Eligible workflows & NDC workflows & NDC rate\\
\hline
npm-compatible only & 14,189 & 5,772 & 40.68\%\\
Python-compatible only & 13,419 & 9,707 & 72.34\%\\
Both groups & 842 & 659 & 78.27\%\\
\hline
Total & 28,450 & 16,138 & 56.72\%\\
\hline
\end{tabular}
\end{table}

\section*{Supplementary repository-feature distributions}

Supplementary Fig.~S1 shows workflow prevalence by repository-feature stratum. Supplementary
Figs.~S2--S12 show the complete set of 11 CCS-specific repository-feature
comparisons. In every split violin, the gray left half is the
same global reference of repositories with no detectable CCS, whereas the
blue right half contains repositories with the named target CCS. Target
repositories may contain other CCS types. These global-reference comparisons are
unadjusted, are not detector-opportunity matched, and do not estimate
independent or causal effects.

\begin{figure}[H]
\centering
\includegraphics[width=\linewidth]{figures/repository_feature_smell_heatmap_canonical_rq1_20260921.png}
\caption{Unadjusted workflow prevalence of the 11 statically detectable CCS within repository-feature strata. Each cell uses its stratum-specific workflow denominator; missing values are retained in source data and omitted from the display. The heatmap and continuous-distribution comparisons use the same final RQ1 detector results.}
\end{figure}

\clearpage
\begin{figure}[H]
\centering
\includegraphics[width=\linewidth]{figures/no_dependency_cache_global_nosmell_comparison_ccs.pdf}
\caption{Repository characteristics for \texttt{no\_dependency\_cache}. Each panel compares repositories with no detectable CCS (left half, gray; $n=8{,}796$, except contributors $n=8{,}794$) with repositories containing the target CCS (right half, blue; $n=6{,}976$). Features use the $\log_{10}(x+1)$ scale; black bars and points show the interquartile range and median.}
\end{figure}

\clearpage
\begin{figure}[H]
\centering
\includegraphics[width=\linewidth]{figures/no_fallback_cache_global_nosmell_comparison_ccs.pdf}
\caption{Repository characteristics for \texttt{no\_fallback\_cache}. Each panel compares repositories with no detectable CCS (left half, gray; $n=8{,}796$, except contributors $n=8{,}794$) with repositories containing the target CCS (right half, blue; $n=3{,}007$). Features use the $\log_{10}(x+1)$ scale; black bars and points show the interquartile range and median.}
\end{figure}

\clearpage
\begin{figure}[H]
\centering
\includegraphics[width=\linewidth]{figures/redundant_cache_update_global_nosmell_comparison_ccs.pdf}
\caption{Repository characteristics for \texttt{redundant\_cache\_update}. Each panel compares repositories with no detectable CCS (left half, gray; $n=8{,}796$, except contributors $n=8{,}794$) with repositories containing the target CCS (right half, blue; $n=1{,}460$). Features use the $\log_{10}(x+1)$ scale; black bars and points show the interquartile range and median.}
\end{figure}

\clearpage
\begin{figure}[H]
\centering
\includegraphics[width=\linewidth]{figures/cache_position_error_global_nosmell_comparison_ccs.pdf}
\caption{Repository characteristics for \texttt{cache\_position\_error}. Each panel compares repositories with no detectable CCS (left half, gray; $n=8{,}796$, except contributors $n=8{,}794$) with repositories containing the target CCS (right half, blue; $n=528$). Features use the $\log_{10}(x+1)$ scale; black bars and points show the interquartile range and median.}
\end{figure}

\clearpage
\begin{figure}[H]
\centering
\includegraphics[width=\linewidth]{figures/overlapping_cache_mechanisms_global_nosmell_comparison_ccs.pdf}
\caption{Repository characteristics for \texttt{overlapping\_cache\_mechanisms}. Each panel compares repositories with no detectable CCS (left half, gray; $n=8{,}796$, except contributors $n=8{,}794$) with repositories containing the target CCS (right half, blue; $n=407$). Features use the $\log_{10}(x+1)$ scale; black bars and points show the interquartile range and median.}
\end{figure}

\clearpage
\begin{figure}[H]
\centering
\includegraphics[width=\linewidth]{figures/static_cache_key_global_nosmell_comparison_ccs.pdf}
\caption{Repository characteristics for \texttt{static\_cache\_key}. Each panel compares repositories with no detectable CCS (left half, gray; $n=8{,}796$, except contributors $n=8{,}794$) with repositories containing the target CCS (right half, blue; $n=1{,}142$). Features use the $\log_{10}(x+1)$ scale; black bars and points show the interquartile range and median.}
\end{figure}

\clearpage
\begin{figure}[H]
\centering
\includegraphics[width=\linewidth]{figures/overly_specific_cache_key_global_nosmell_comparison_ccs.pdf}
\caption{Repository characteristics for \texttt{overly\_specific\_cache\_key}. Each panel compares repositories with no detectable CCS (left half, gray; $n=8{,}796$, except contributors $n=8{,}794$) with repositories containing the target CCS (right half, blue; $n=405$, except contributors $n=404$). Features use the $\log_{10}(x+1)$ scale; black bars and points show the interquartile range and median.}
\end{figure}

\clearpage
\begin{figure}[H]
\centering
\includegraphics[width=\linewidth]{figures/no_docker_layer_cache_global_nosmell_comparison_ccs.pdf}
\caption{Repository characteristics for \texttt{no\_docker\_layer\_cache}. Each panel compares repositories with no detectable CCS (left half, gray; $n=8{,}796$, except contributors $n=8{,}794$) with repositories containing the target CCS (right half, blue; $n=2{,}343$). Features use the $\log_{10}(x+1)$ scale; black bars and points show the interquartile range and median.}
\end{figure}

\clearpage
\begin{figure}[H]
\centering
\includegraphics[width=\linewidth]{figures/deprecated_cache_action_global_nosmell_comparison_ccs.pdf}
\caption{Repository characteristics for \texttt{deprecated\_cache\_action}. Each panel compares repositories with no detectable CCS (left half, gray; $n=8{,}796$, except contributors $n=8{,}794$) with repositories containing the target CCS (right half, blue; $n=249$). Features use the $\log_{10}(x+1)$ scale; black bars and points show the interquartile range and median.}
\end{figure}

\clearpage
\begin{figure}[H]
\centering
\includegraphics[width=\linewidth]{figures/unreachable_cache_write_global_nosmell_comparison_ccs.pdf}
\caption{Repository characteristics for \texttt{unreachable\_cache\_write}. Each panel compares repositories with no detectable CCS (left half, gray; $n=8{,}796$, except contributors $n=8{,}794$) with repositories containing the target CCS (right half, blue; $n=4$). The sparse target group is marked $\dagger$, and all target observations are overlaid; its kernel-density shape is unstable and should not be generalized. Features use the $\log_{10}(x+1)$ scale; black bars and points show the interquartile range and median.}
\end{figure}

\clearpage
\begin{figure}[H]
\centering
\includegraphics[width=\linewidth]{figures/cache_as_artifact_global_nosmell_comparison_ccs.pdf}
\caption{Repository characteristics for \texttt{cache\_as\_artifact}. Each panel compares repositories with no detectable CCS (left half, gray; $n=8{,}796$, except contributors $n=8{,}794$) with repositories containing the target CCS (right half, blue; $n=116$). The sparse target group is marked $\dagger$, and all target observations are overlaid; its kernel-density shape should be interpreted cautiously. Features use the $\log_{10}(x+1)$ scale; black bars and points show the interquartile range and median.}
\end{figure}

\end{document}
